\documentclass[11pt]{article}
\usepackage{fontspec}
\setmainfont{Times New Roman}
\setsansfont{Arial}
\setmonofont{Consolas}
\usepackage{amsmath,amssymb}
\usepackage{geometry}
\usepackage{microtype}
\geometry{a4paper,margin=3cm}
\setlength{\parindent}{1.25em}
\setlength{\parskip}{0pt}
\emergencystretch=10em
\allowdisplaybreaks
\providecommand{\Kfield}{\mathfrak{K}}
\providecommand{\Ssys}{\mathfrak{S}}
\providecommand{\Jdom}{\mathfrak{J}}
\providecommand{\Gdom}{\mathfrak{G}}
\providecommand{\Lsys}{\mathfrak{L}}
\providecommand{\Omegaint}{[\Omega]}
\providecommand{\eps}{\varepsilon}

\begin{document}

\section*{\S\ 10. Релативно целе области првого рода и ньихова интегрална база.}

Појем ``релативно алгебраично целого'', введены в \S\ 9 (Дефиниција В), даваје можност указати класу релативно целых областиј, кторе сут конечне интегралне области, и тым самым најманье за ту класу позитивно одговорити на пытанье, поставјено од D. Хилберта (Матхематисцхе Проблеме, проблема 14). Дајемо следујучу

{\itshape Дефиниција VI: Конечно число полиномов из $\Gdom_{n\rho}$ зове се интегрална база, ако всак полином из $\Gdom_{n\rho}$ можно представити како целу рационалну комбинацију того конечного числа с коефициентами из $\Omega$. Интегрална област из полиномов, ктора имаје интегралну базу, зове се конечна интегрална област.}

Класу релативно целых областиј, ктору разматрајемо и за ктору се инволуцијска база покаже како интегрална база, будемо называти релативно целыми областјами првого рода по следујучеј дефиницији:

{\itshape Дефиниција VII: Релативно цела област $\Gdom_{n\rho}$ зове се област првого рода, ако како област одображеньа имаје релативно целу област $\Gdom^*_{\rho\rho}$ таку, же всак либовольны полином од $x_1,\ldots,x_\rho$ (с коефициентами из $\Omega$) јест алгебраично целы над $\Gdom^*_{\rho\rho}$.}

По тој дефиницији неопредельене $x_1,\ldots,x_\rho$, и следователно такоже линеарна форма $u(x)=u_1x_1+\cdots+u_\rho x_\rho$, сут алгебраично целе над $\Gdom^*_{\rho\rho}$. Зато неразложима цела функција с кореном $z=u(x)$ --- инволуцијска форма најменшего польа $\Kfield^*_{\rho\rho}$, содержечего $\Gdom^*_{\rho\rho}$ --- имаје како највышши коефициент јединицу, а како дальне коефициенты полиномы из $\Gdom^*_{\rho\rho}$. Зато она једночасно јест инволуцијска форма $\Gdom^*_{\rho\rho}$, и в $\Gdom^*_{\rho\rho}$ не може ексистовати друга инволуцијска форма. По принципу преноса $\Gdom_{n\rho}$ такоже имаје инволуцијску форму, чији највышши коефициент јест јединица; а понеже по Теореме VI в рационалном представјеньу всех полиномов из $\Gdom_{n\rho}$ чрез соответну инволуцијску базу в именователь входи само тој највышши коефициент, инволуцијска база ставаје интегрална база; то јест:

{\itshape Теорема VII. Всака релативно цела област првого рода јест конечна интегрална област; интегрална база дана јест инволуцијскоју базоју, добытоју из характеризујучеј области одображеньа. Особливо, интегрална база релативно целых областиј првого рода $\Gdom_{\rho\rho}$ дана јест једнозначно дефинованоју инволуцијскоју базоју $\Gdom_{\rho\rho}$.}

Дајемо јешче критериј за релативно целе области првого рода. Нехај $\Gdom_{n\rho}$ јест област првого рода, и нехај интегрална база $\Gdom_{n\rho}$ јест дана полиномами
\begin{equation}
  F_1(x),F_2(x),\ldots,F_k(x).
\tag{1}
\end{equation}
Тогда по Хилбертовом помочном теорему\footnote{D. Хилберт: Üбер дие воллен Инвариантенсыстеме, Math. Анн. 42 (1893), \S\ 1. (Помочны теорем, там изречены само за хомогене полиномы, једначно важи за нехомогене.)} ексистујут $\rho$ алгебраично независне полиномы из $\Gdom_{n\rho}$:
\begin{equation}
  G_1(x),G_2(x),\ldots,G_\rho(x)
\tag{2}
\end{equation}
таке, же полиномы (1), и зато всак полином из $\Gdom_{n\rho}$, сут алгебраично целе над полиномами (2). То само важи за всак полином области одображеньа $\Gdom^*_{\rho\rho}$ взходно к соответным полиномам
\begin{equation}
  G_1^*(x),G_2^*(x),\ldots,G_\rho^*(x).
\tag{3}
\end{equation}

По Дефиницији VII зато такоже всак либовольны полином од $x_1,\ldots,x_\rho$ јест алгебраично целы над полиномами (3).

Обратно, нехај област одображеньа $\Jdom^*_{\rho\rho}$ либовольној релативно целој области $\Gdom_{n\rho}$ содержи $\rho$ полиномов (3), над кторыми всак полином од $x_1,\ldots,x_\rho$ јест алгебраично целы. Тогда покажемо, же из того $\Jdom^*_{\rho\rho}$ сама настава како релативно цела област $\Gdom^*_{\rho\rho}$, и далье, же $\Gdom^*_{\rho\rho}$, а следователно и $\Gdom_{n\rho}$, сут првого рода.

В самој ствари, нехај $F^*(x)$ јест полином из најменшего польа $\Kfield^*_{\rho\rho}$, содержечего $\Jdom^*_{\rho\rho}$; по предпоставце он јест алгебраично целы над полиномами (3). Тогда соответнаја функција из $\Kfield_{n\rho}$ јест алгебраично цела над $\rho$ полиномами из $\Gdom_{n\rho}$, зато јест полином $F(x)$ и следователно належи к $\Gdom_{n\rho}$; с тым $F^*(x)$ належи к $\Jdom^*_{\rho\rho}$. Област одображеньа $\Jdom^*_{\rho\rho}$ содержи зато все полиномы из $\Kfield^*_{\rho\rho}$ и јест релативно цела област $\Gdom^*_{\rho\rho}$. За $\Gdom^*_{\rho\rho}$, и зато такоже за $\Gdom_{n\rho}$, из предпоставкы по Дефиницијах В и VII непосреднье следује, же оны сут првого рода, то јест:

{\itshape Релативно целе области првого рода $\Gdom_{n\rho}$ характеризоване сут тым, же в области одображеньа $\Jdom^*_{\rho\rho}$ ексистујут $\rho$ полиномов, над кторыми всак либовольны полином од $x_1,\ldots,x_\rho$ јест алгебраично целы. Из тој предпоставкы област одображеньа сама настава како релативно цела област.}\footnote{Например, в области $\Gdom_{22}$, споменутој в конечном замеченьу к \S\ 9, допустима специализација $x_1=1$ даје два полиномы $F^*=x_2$, $G^*=1+x_3$, над кторыми всак либовольны полином од $x_2,x_3$ јест алгебраично целы. Зато $\Gdom_{22}$ јест првого рода, именно с $F(x)$ и $G(x)$ како интегралноју базоју.}

\end{document}

